\scanpage{7}
$X\xrightarrow{f}S$ morphisme \deleted{régulier} propre

$X$ régulier

$S$ régulier. irréductible, pt. gén. $\eta$

On suppose que $\underline{\operatorname{Pic}}_{X/S}^{\deleted{\unsure{0}}}$ existe, et \deleted{est une V.A.}
\deleted{Par} que l’adhérence \inserted{schématique} de $A_\eta=\underline{\operatorname{Pic}}_{X_\eta/k(\eta)}^{0}$, \deleted{\unsure{est}}\unsure{soit} un
\uline{sous-schéma en groupes} \uline{plat} de $X$, \deleted{ce qui}
\uline{lisse} sur $S$ \deleted{\unsure{a|\textemdash}} ce qui est automatiquement
le cas si $\underline{\operatorname{Pic}}^{\unsure{0}}_{X/S}$ est lisse sur $S$ le long de
la section unité, p.ex. si $S$ régulier de dim $1$
et les $\operatorname{Pic}_{X_s/k(s)}$ sont lisses $s\in S$).

Alors \deleted{$\operatorname{Pic}$} $A$ \uline{est un modèle de Néron} :
en effet, si $S'/S$ est lisse, alors
$X'=X\times_S S'$ est lisse sur $X$ régulier, donc est
régulier, donc tt \deleted{\unsure{sectio|section}} élément de $\operatorname{Pic}(X'_{\unsure{\eta'|\eta}})$
provient d’un élément de $\operatorname{Pic}(X')$.

% slanted
\margin{Attention s’il n’y a pas de section locale \inserted{(ét.)} de $X$ sur $S$}
% archive sheet number: 6
